% !TEX root = ../matan.tex

\section{Ортонормированный базис в сепарабельном гильбертовом пространстве.}

\begin{Definition}{Сепарабельность}{}
	Гильбертово пространство $\mathcal{H}$ называется \textbf{сепарабельным}, если в нём существует счётное всюду плотное подмножество, то есть
	\[
	\exists\, E = \{y_k\}_{k \in \NN} \subset \mathcal{H} \st \operatorname{clos} E = \mathcal{H}.
	\]
	Эквивалентно: $\forall x \in \mathcal{H}\ \exists\,\{x_k\} \subset E \st \|x_k - x\|_{\mathcal{H}} \to 0$.
\end{Definition}

Напомним, что \textbf{ортонормированный базис}~--- это полная ортонормированная система: $\langle e_i, e_j\rangle = \delta_{ij}$ и $\operatorname{clos}\operatorname{span}\{e_k\} = \mathcal{H}$.

\begin{Theorem}{Счётность ортонормированной системы}{}
	В сепарабельном гильбертовом пространстве $\mathcal{H}$ всякая ортонормированная (а значит, и всякая ортогональная) система не более чем счётна.
\end{Theorem}

\begin{proof}
	Достаточно рассмотреть ортонормированную систему $\{x_\alpha\}_{\alpha \in A}$ (общий ортогональный случай сводится к этому нормировкой, не меняющей мощности системы). Для разных $\alpha \neq \beta$
	\[
	\|x_\alpha - x_\beta\|
	= \sqrt{\langle x_\alpha, x_\alpha\rangle - 2\langle x_\alpha, x_\beta\rangle + \langle x_\beta, x_\beta\rangle}
	= \sqrt{1 - 0 + 1} = \sqrt{2}.
	\]
	Рассмотрим открытые шары $B_\alpha = \bigl\{ x \st \|x - x_\alpha\| < \tfrac12 \bigr\}$. Они попарно не пересекаются: если бы $z \in B_\alpha \cap B_\beta$, то по неравенству треугольника
	\[
	\sqrt{2} = \|x_\alpha - x_\beta\| \le \|x_\alpha - z\| + \|z - x_\beta\| < \tfrac12 + \tfrac12 = 1,
	\]
	что невозможно. Значит, $B_\alpha \cap B_\beta = \varnothing$ при $\alpha \neq \beta$.

	Так как $\mathcal{H}$ сепарабельно, существует счётное всюду плотное $E = \{y_k\}_{k \in \NN}$. По плотности каждый шар $B_\alpha$ содержит хотя бы одну точку $y_{n(\alpha)} \in E \cap B_\alpha$. Поскольку шары не пересекаются, отображение $\alpha \mapsto n(\alpha)$ инъективно: разным $\alpha$ отвечают точки из разных непересекающихся шаров, а значит, разные элементы $E$. % [восстановлено]
	Итак, индексное множество $A$ инъективно вкладывается в счётное множество $\NN$, поэтому $A$ не более чем счётно.
\end{proof}

\begin{Theorem}{Существование ортонормированного базиса}{}
	В любом сепарабельном гильбертовом пространстве $\mathcal{H} \neq \{0\}$ существует не более чем счётный ортонормированный базис.
\end{Theorem}

\begin{proof}
	Возьмём счётное всюду плотное множество $E = \{y_k\}_{k \in \NN}$. Выделим из него линейно независимую подсистему $\{h_m\}$ с той же линейной оболочкой: идём по $k = 1, 2, \ldots$ и оставляем очередной $y_k$ тогда и только тогда, когда он не лежит в линейной оболочке уже отобранных. Тогда $\operatorname{span}\{h_m\} = \operatorname{span}\{y_k\}$, поэтому
	\[
	\operatorname{clos}\operatorname{span}\{h_m\} = \operatorname{clos}\operatorname{span}\{y_k\} \supseteq \operatorname{clos} E = \mathcal{H},
	\]
	то есть система $\{h_m\}$ линейно независима и полна.

	Применим к $\{h_m\}$ процесс ортогонализации Грама--Шмидта (см. вопрос~18): получим ортонормированную систему $\{e_m\}$ с
	\[
	\operatorname{span}\{e_1, \ldots, e_m\} = \operatorname{span}\{h_1, \ldots, h_m\} \qquad \forall m.
	\]
	Отсюда $\operatorname{span}\{e_m\} = \operatorname{span}\{h_m\}$, а значит, $\operatorname{clos}\operatorname{span}\{e_m\} = \mathcal{H}$. Таким образом, $\{e_m\}$~--- полная ортонормированная система, то есть ортонормированный базис; его счётность гарантирует предыдущая теорема.
\end{proof}

\begin{Remark}{}{}
	Существование такого базиса позволяет каждому $x \in \mathcal{H}$ сопоставить набор коэффициентов Фурье $c_k = \langle x, e_k\rangle$ (вопросы~19--20); именно базис делает это сопоставление взаимно однозначным и изометричным.
\end{Remark}
